\documentclass[11pt]{article}

\usepackage{fontspec}
\setmainfont{TeX Gyre Pagella}
% Container lacks the Pagella Math OTF; using classic math setup here.
% For local builds, restore:
\usepackage{unicode-math}
\setmathfont{TeX Gyre Pagella Math}
\usepackage{amsmath,amssymb}

\usepackage[margin=1.1in]{geometry}
\usepackage{amsthm}
\usepackage{booktabs}
\usepackage{microtype}
\emergencystretch=1.5em
\usepackage{tcolorbox}

\newtheorem{theorem}{Theorem}[section]
\newtheorem{proposition}[theorem]{Proposition}
\newtheorem{lemma}[theorem]{Lemma}
\newtheorem{conjecture}[theorem]{Conjecture}
\newtheorem{corollary}[theorem]{Corollary}
\newtheorem{definition}[theorem]{Definition}
\newtheorem{problem}[theorem]{Problem}
\theoremstyle{remark}
\newtheorem{remark}[theorem]{Remark}

\newcommand{\C}{\mathcal{C}}
\newcommand{\Ops}{\mathcal{O}}
\newcommand{\A}{\mathcal{A}}
\newcommand{\R}{\mathsf{R}}
\newcommand{\E}{\mathsf{E}}
\newcommand{\G}{\mathsf{G}}
\newcommand{\Tor}{\operatorname{Tor}}

\title{Repair Categories and the Local Geometry of Continuation}
\author{Flyxion}
\date{July 2026 \textperiodcentered\ Draft 0.1}

\begin{document}
\maketitle

\begin{abstract}
A continuation system carries two independent notions of proximity: a
directed \emph{repair contrast} $\R(C,C')$, measuring the cost of
recognizing $C'$ as a repaired form of $C$, and an \emph{executable
cost} $\E(C,C')$, the least total cost of an operational path from $C$
to $C'$. We study the geometry these jointly induce, beginning from a
fully explicit twelve-state system in which every quantity is
computable by hand. The computation isolates a fork that was not
visible at the level of analogy: if $\R$ satisfies the triangle
inequality --- that is, if repair space is a Lawvere-enriched category
--- then the Pythagorean projection defect onto any admissible subset
is necessarily nonpositive, and coherence between comparison and
execution is certifiable by a local check on generators. Positive
projection interference, the signature phenomenon of dually flat
information geometry, requires a divergence-type contrast and admits
no local coherence certificate. We then pose and solve the finite
characterization problem in the style of Chentsov. Monotonicity alone
is maximally underdetermined: natural monotone invariants of the
category separate every ordered pair of the model, detecting
operational data --- promotion versus phase, proximity to the
singular stratum --- that the enrichment forgets. Adding invariance
under repair-sufficient morphisms collapses them exactly:
sufficiency, provably scarce inside irreversible systems, can be
manufactured between systems by amalgamation, and the invariants so
characterized are precisely the monotone functions of the
\emph{bidirected} executable cost, the ordered pair of forward and
reverse repair cost. Directedness survives characterization;
operational-bracket data survives only monotone naturality. The
smooth dual-geometric theory is presented as the tangent
approximation to the enriched substrate, not its foundation.
\end{abstract}

\section*{Introduction}

Most theories of continuation assume a fixed structure of comparison.
An object changes; the framework for measuring change stands still.
The work collected under the Admissibility Program has repeatedly
encountered the opposite situation: admissibility conditions vary
across a space of contexts, repair criteria drift, projection systems
reorganize, and retained histories differ from point to point. In
such a setting the naive difference between a continuation here and a
continuation there conflates two things --- change in the object and
change in the frame --- and the problem of separating them is,
structurally, the problem that led differential geometry to the
connection and the covariant derivative. The tempting conclusion is
that a theory of admissibility should import that machinery. This
paper argues, and in the finite case proves, that the inheritance
runs the other way. The differential-geometric apparatus is a tangent
approximation to something more primitive, procedural and
enriched-categorical in kind; and when one asks --- in the manner of
Chentsov --- which quantities are actually \emph{characterized} by
the invariances a continuation system supports, the answer is not a
connection at all.

Every continuation system carries two proximities, and they are
this paper's primitives. The \emph{executable cost} $\E(C,C')$ is the least
total cost of a sequence of available operations transforming $C$
into $C'$; it is generated by the system's own operational algebra,
and its triangle inequality is nothing but the composition of
procedures --- in Lawvere's sense, $\E$ is the free enrichment of the
operational graph. The \emph{repair contrast} $\R(C,C')$ is a
directed comparison defined independently of the operations: the cost
of recognizing $C'$ as a repaired form of $C$. Nothing forces these
to agree. Some states are easy to compare but difficult to reach;
some are easy to reach but expensive to repair. The discrepancy
$\G = \E - \R$ between comparison and execution is the central
object of study.

Everything in the paper answers to a single worked example: a
twelve-state system with two operations, one of them irreversible,
small enough that every claim reduces to a table. The example was
built to discipline the abstractions, and it did more than that; each
of the paper's results was found in its tables before it was proved.
Three findings organize the paper. First, a \emph{fork}: if $\R$
satisfies the triangle inequality --- if repair space is itself an
enriched category --- then the Pythagorean projection defect onto any
admissible subset is necessarily nonpositive, and coherence between
comparison and execution ($\R\le\E$) is certifiable by a local check
on generators. Positive projection interference, the signature of
dually flat information geometry, is possible only for
divergence-type contrasts, which admit no such local certificate.
The enriched and divergence regimes of repair are separated by a
computable sign. Second, a \emph{profusion}: posing the finite
characterization problem --- which functionals of pairs of states are
monotone under all admissibility morphisms --- the answer is that
monotonicity alone constrains almost nothing. The natural monotone
invariants separate every ordered pair of the model, and what they
detect is operational data the enrichment forgets: which generator
does the work of a geodesic, and how close it runs to the singular
stratum where the operations cease to commute. Third, a
\emph{collapse}, which is the paper's main theorem: adding invariance
under repair-sufficient morphisms --- the analogue of Chentsov's
sufficient statistics --- cuts the profusion down exactly. Sufficient
morphisms are provably scarce inside an irreversible system, but they
can be \emph{manufactured} between systems by amalgamation, gluing
two systems along a pair of states with matching costs; and the
invariants that survive are precisely the monotone functions of the
\emph{bidirected} executable cost, the ordered pair
$\bigl(\E(x,y),\E(y,x)\bigr)$. What survives every
distinction-preserving redescription of a continuation system is two
numbers per ordered pair: the cost of getting there, and the cost of
getting back.

The theorem stands to Chentsov's as irreversible systems stand to
statistical ones, with the hinge inverted: where the classical proof
leans on the abundance of sufficient statistics, the present proof
leans on a construction, because abundance fails. And the
characterized invariant is not a symmetric metric but a directed
pair; directedness survives characterization, while the
operational-bracket data --- the torsion-like, order-sensitive
structure --- survives only the weaker discipline of monotone
naturality. The differential-geometric objects that motivated the
inquiry reappear at the end, in their proper place: on the simplex
relaxation of a finite system, the metric and dual connections of
information geometry arise as the second- and third-order jets of a
smooth contrast, and the paper's closing conjecture is that the
smooth characterization, when it exists, will pin the germ of the
bidirected cost --- the dual geometry emerging as the jet expansion
of the characterized invariant rather than as a structure selected by
fiat.

Section~1 presents the twelve-state system and the fork. Section~2
develops the enriched-categorical foundation. Section~3 poses and
solves the finite characterization problem. Section~4 constructs the
smooth tangent theory and states what remains open.

\section{A Twelve-State Continuation System}

We begin from something almost embarrassingly concrete: an example
small enough that every claim in this paper reduces to a table.

\begin{definition}[The system]
Let $\C = \{0,1,2\}\times\{0,1,2,3\}$, a state $C=(\ell,p)$ consisting
of a \emph{level} $\ell$ and a cyclic \emph{phase} $p$. The
operational algebra $\Ops$ is generated by two total maps:
\[
a(\ell,p) = (\ell,\, p+1 \bmod 4), \qquad
b(\ell,p) =
\begin{cases}
(\ell+1,\,p) & \ell < 2,\\
(2,\,0) & \ell = 2,
\end{cases}
\]
with state-independent costs $c(a)=1$ and $c(b)=2$. No operation
decreases $\ell$: promotion is irreversible, and saturation of $b$ at
the top level resets the phase.
\end{definition}

\begin{definition}[Executable cost]
$\E(C,C')$ is the infimum of $\sum_i c(o_i)$ over finite words
$o_n\cdots o_1$ in $\{a,b\}$ with $o_n\cdots o_1(C)=C'$, and $+\infty$
if no such word exists.
\end{definition}

Metric geometry begins from distance. Repair geometry begins from
reconstruction.

\begin{definition}[Repair contrast]
Independently of $\Ops$, define the directed contrast
\[
\R\bigl((\ell,p),(\ell',p')\bigr)
= \max(\ell'-\ell,0) + 5\max(\ell-\ell',0)
+ \tfrac12\, d_{\mathrm{cyc}}(p,p'),
\]
where $d_{\mathrm{cyc}}$ is cyclic distance on $\{0,1,2,3\}$.
Recognizing an upgrade is cheap; recognizing a downgrade is expensive
but finite; phase mismatch costs $\tfrac12$ per cyclic step. Note that
$\R$ satisfies the triangle inequality: it is a quasi-metric, hence a
Lawvere enrichment of $\C$.
\end{definition}

Three quantities are computed exhaustively over all $132$ ordered
pairs (source and full computation in the appendix script).

\paragraph{The comparison--execution gap.}
Define $\G = \E - \R \in [0,\infty]$. Of the $132$ pairs, $48$ have
$\G=\infty$: every downward repair is \emph{recognizable} at finite
contrast ($5$ per level) but \emph{inexecutable}, since no operation
lowers $\ell$. Among finite pairs the largest gaps occur where $\R$
exploits cyclic symmetry that $\Ops$ cannot: for instance
$\R\bigl((0,0),(2,3)\bigr)=2.5$ while $\E\bigl((0,0),(2,3)\bigr)=7$,
because the phase must be cranked forward three times and two
promotions paid. Representational nearness and operational
accessibility dissociate, and the dissociation is quantitative.

\begin{lemma}[Local-to-global coherence]\label{lem:coherence}
Suppose $\R$ satisfies the triangle inequality and is
\emph{generator-dominated}: $\R(C,o(C)) \le c(o;C)$ for every
generator $o\in\Ops$ and every $C\in\C$. Then $\R \le \E$, i.e.\
$\G\ge0$, on all of $\C\times\C$.
\end{lemma}

\begin{proof}
For any executable word $o_n\cdots o_1$ from $C$ to $C'$ with
intermediate states $C_i$, the triangle inequality gives
$\R(C,C') \le \sum_i \R(C_{i-1},C_i) \le \sum_i c(o_i;C_{i-1})$.
Take the infimum over words.
\end{proof}

In the model, generator domination holds with room to spare ($a$ moves
$\R$ by at most $\tfrac12 \le 1$; $b$ by at most $1 \le 2$), so
$\G\ge0$ is a theorem of the system rather than an empirical
observation. The lemma fails without the triangle inequality; we
return to this below.

\paragraph{Localized operational torsion.}
The generators commute everywhere except on the saturation stratum:
for every $p$,
\[
a\circ b\,(2,p) = (2,1), \qquad b\circ a\,(2,p) = (2,0),
\]
while $a\circ b = b\circ a$ on $\ell<2$. Order-dependence is not
diffuse; it is concentrated exactly where $b$ changes character from
injective promotion to non-injective collapse.

\begin{proposition}[Localized torsion witness]
The commutator defect of\/ $\Ops$ is supported on the singular stratum
$\{\ell=2\}$ of the operational algebra.
\end{proposition}

\paragraph{The projection defect and its sign.}
For an admissible subset $\A\subseteq\C$ with projection
$C^{*} = \arg\min_{t\in\A}\R(C,t)$, define
\[
\Delta_{\A}(C,C') = \R(C,C') - \R(C,C^{*}) - \R(C^{*},C'),
\qquad C'\in\A .
\]
With the level-aligned set $\A=\{(1,p): p\}$ we find
$\Delta_{\A}\equiv 0$: the contrast is separable and $\A$ is
autoparallel with respect to the product structure. With the diagonal
set $\A=\{(0,0),(1,2),(2,1)\}$, fourteen of thirty-six pairs give
$\Delta_{\A}\neq 0$, with range $[-2,0]$. The defect never becomes
positive, and this is not an accident of the example.

\begin{proposition}[Sign restriction in enriched repair]
\label{prop:sign}
If $\R$ satisfies the triangle inequality, then
$\Delta_{\A}(C,C')\le 0$ for every $\A$, every projection point, and
every $C'\in\A$. Consequently
\[
\Delta_{\A} > 0 \ \Longrightarrow\ \R \text{ is not a Lawvere
enrichment.}
\]
\end{proposition}

\begin{proof}
Immediate from $\R(C,C')\le \R(C,C^{*})+\R(C^{*},C')$.
\end{proof}

\begin{tcolorbox}[title={The fork}]
Positive Pythagorean interference --- the signature of dually flat
information geometry, where the generalized Pythagorean theorem yields
exact positive decompositions --- is \emph{impossible} in enriched
repair. At the same location, Lemma~\ref{lem:coherence} shows that
enriched repair admits a \emph{local certificate} of coherence
($\R\le\E$ verified generator by generator), whereas a divergence-type
contrast, lacking the triangle inequality, admits none: its coherence
with execution must be imposed globally. Repair structures therefore
bifurcate into an \emph{enriched regime} (directed metric, locally
certifiable, projection can only lose) and a \emph{divergence regime}
(contrast function, dual connections, positive interference possible,
coherence a global axiom). The remainder of this paper develops the
first regime rigorously and situates the second as its tangent theory.
\end{tcolorbox}

\section{Repair as a Lawvere-Enriched Category}

The twelve-state system used no structure beyond a set of states, a
costed operational graph, and a directed contrast. This section fixes
that data abstractly and shows that the executable cost $\E$ is not an
auxiliary construction but the free object of the theory. Throughout,
costs are valued in the quantale
$\mathbb{V} = ([0,\infty],\,\ge,\,+)$, the base of enrichment in
Lawvere's generalized metric spaces \cite{lawvere}: objects of a
$\mathbb{V}$-category are points, the hom from $x$ to $y$ is a cost
$d(x,y)\in[0,\infty]$ with $d(x,x)=0$, and the composition law
\emph{is} the triangle inequality
$d(x,z) \le d(x,y) + d(y,z)$.

\begin{definition}[Continuation system]\label{def:cs}
A \emph{finite continuation system} is a triple $(\C,\Ops,c)$ where
$\C$ is a finite set of states, $\Ops$ is a set of generating
operations $o\colon \operatorname{dom}(o)\subseteq\C \to \C$, and
$c(o;C)\in(0,\infty)$ assigns each application its cost. The
\emph{operational graph} has vertex set $\C$ and a directed edge
$C \to o(C)$ of weight $c(o;C)$ for each applicable generator.
\end{definition}

\begin{definition}[Executable enrichment]\label{def:free}
The \emph{executable cost} $\E(C,C')$ is the infimum of
$\sum_{i} c(o_i;C_{i-1})$ over finite operational words from $C$ to
$C'$, with $\E(C,C)=0$ (the empty word) and $\E(C,C')=\infty$ when no
word exists.
\end{definition}

$\E$ satisfies the triangle inequality by concatenation of words, so
$(\C,\E)$ is a $\mathbb{V}$-category. It is, moreover, the
\emph{free} one on the operational graph: enrichment is not an
assumption about repair but a consequence of composition.

\begin{definition}[Admissibility morphism]\label{def:morphism}
A map $F\colon (\C_1,\Ops_1,c_1) \to (\C_2,\Ops_2,c_2)$ between
continuation systems is an \emph{admissibility morphism} if there is
an assignment $o \mapsto F_{*}(o)$ of generators to (possibly empty)
operational words such that
\[
F\bigl(o(C)\bigr) = F_{*}(o)\bigl(F(C)\bigr)
\quad\text{and}\quad
c_2\bigl(F_{*}(o); F(C)\bigr) \le c_1(o;C)
\]
for every applicable pair, where $c_2$ of a word is the sum of its
step costs. Intertwining says $F$ respects what can be done;
cost-nonincrease says $F$ never makes an operation more expensive.
Coarse-grainings, quotients by zero-cost equivalences, and forgetful
projections are all of this form.
\end{definition}

\begin{lemma}[Nonexpansiveness]\label{lem:nonexp}
Every admissibility morphism satisfies
$\E_2\bigl(F(C),F(C')\bigr) \le \E_1(C,C')$.
\end{lemma}

\begin{proof}
The $F_{*}$-image of an operational word from $C$ to $C'$ is an
operational word from $F(C)$ to $F(C')$ of no greater total cost;
take infima.
\end{proof}

In enriched language, Lemma~\ref{lem:nonexp} says every admissibility
morphism is a $\mathbb{V}$-functor $(\C_1,\E_1)\to(\C_2,\E_2)$. The
next proposition locates $\E$ uniquely among all contrasts compatible
with the generators; it is the abstract form of
Lemma~\ref{lem:coherence}.

\begin{proposition}[Maximality of the free enrichment]
\label{prop:maximal}
Let $d\colon \C\times\C \to [0,\infty]$ satisfy $d(C,C)=0$, the
triangle inequality, and generator domination
$d(C,o(C)) \le c(o;C)$. Then $d \le \E$ pointwise. Hence $\E$ is the
greatest generator-dominated enrichment on $(\C,\Ops,c)$, and
$\G = \E - \R \ge 0$ holds for \emph{every} enriched, generator-dominated
repair contrast $\R$ --- Lemma~\ref{lem:coherence} is its instance.
\end{proposition}

\begin{proof}
Chain $d$ along any word and take infima, as in
Lemma~\ref{lem:coherence}. Freeness in the categorical sense: $\E$ is
the left Kan extension of the weighted operational graph along the
inclusion of generators into words, i.e.\ the value of the left
adjoint to the forgetful functor from $\mathbb{V}$-categories to
costed directed graphs.
\end{proof}

Two consequences follow. First, the comparison--execution
gap $\G$ acquires an intrinsic meaning: for enriched contrasts it
measures the failure of $\R$ to be \emph{free}, i.e.\ how far the
system's abstract sense of nearness falls below what its own
operations generate. Second, the sign restriction of
Proposition~\ref{prop:sign} is now visible as a fact about the base of
enrichment: any $\mathbb{V}$-category satisfies $\Delta_{\A}\le 0$,
so positive interference marks a contrast that is not a hom-structure
over $([0,\infty],\ge,+)$ at all.

\begin{definition}[Repair sufficiency]\label{def:sufficient}
An admissibility morphism $F$ is \emph{repair-sufficient} for a
contrast $\R$ if $\R_2\bigl(F(C),F(C')\bigr) = \R_1(C,C')$ for all
$C,C'$.
\end{definition}

Repair-sufficient morphisms coarse-grain without erasing any
repair-relevant distinction; they play the role that sufficient
statistics play in Chentsov's theorem \cite{chentsov} and that
deterministic sufficiency plays in Markov categories \cite{fritz}.
The characterization problem of the next section asks which
functionals are monotone under all admissibility morphisms and
invariant precisely under the repair-sufficient ones.

\begin{remark}[Where the divergence regime lives]
Nothing in Definitions~\ref{def:cs}--\ref{def:sufficient} requires
$\R$ to be triangular. A divergence-type contrast simply fails to be
a $\mathbb{V}$-category structure; it is a two-point function on the
objects of one. The smooth theory of Section~4 will treat such
contrasts as functions on $\Delta(\C)$ whose diagonal jets induce
geometry, with the enriched structure persisting underneath as the
executable skeleton.
\end{remark}

\begin{remark}[Specialization to executable-cost sufficiency]\label{rem:specialize}
Definition~\ref{def:sufficient} defines repair-sufficiency for an
arbitrary contrast $\R$, and nothing said so far ties $\R$ to $\E$.
From this point on, the paper studies sufficiency for the particular
contrast $\R=\E$: condition (I) in Problem~\ref{prob:char} below, and
every theorem that follows from it, characterizes invariance under
\emph{$\E$-sufficient} morphisms, not under morphisms sufficient for
an independently specified $\R$. This is a deliberate narrowing, not
an oversight, but it changes what is being characterized.
Theorem~\ref{thm:main} answers the question ``what survives
redescription that preserves executable cost exactly,'' not ``what
survives redescription that preserves the recognition contrast
$\R$.'' The two questions coincide only when $\R=\E$, which
Definition~\ref{def:cs} does not assume of a general continuation
system. The joint invariant theory of the pair $(\R,\E)$ ---
sufficiency with respect to both primitives simultaneously --- is
left open; it is discussed again in Section~\ref{sec:interp} and in
the closing remarks.
\end{remark}

\section{Characterization in the Finite Case}\label{sec:char}

Chentsov's theorem \cite{chentsov} characterizes the Fisher metric and
the $\alpha$-connections not by construction but by variance: they are
the structures monotone under all Markov morphisms, invariant exactly
under sufficient statistics. The finite analogue for continuation
systems replaces Markov kernels by admissibility morphisms and asks
the corresponding question.

\begin{problem}[Finite characterization]\label{prob:char}
Determine all assignments $\Phi$, giving each finite continuation
system $(\C,\Ops,c)$ a functional
$\Phi_{\C}\colon\C\times\C\to[0,\infty]$, such that
\begin{enumerate}
\item[(M)] $\Phi_{\C_2}\bigl(F(C),F(C')\bigr) \le \Phi_{\C_1}(C,C')$
for every admissibility morphism $F$, and
\item[(I)] equality holds in (M) whenever $F$ is repair-sufficient for
$\E$ (Definition~\ref{def:sufficient} with $\R=\E$).
\end{enumerate}
\end{problem}

Two partial results bound the answer from below and constrain it from
above.

\begin{proposition}[The executable cone]\label{prop:cone}
For every nondecreasing $f\colon[0,\infty]\to[0,\infty]$ with
$f(0)=0$, the assignment $\Phi = f\circ\E$ satisfies (M) and (I). The
class of solutions is closed under pointwise suprema, infima, and
nondecreasing reparametrization, and therefore contains the closed
cone generated by $\E$.
\end{proposition}

\begin{proof}
(M) is Lemma~\ref{lem:nonexp} composed with monotone $f$; (I) is
immediate from the equality defining $\E$-sufficiency. Closure
properties follow pointwise.
\end{proof}

\begin{proposition}[Rigidity]\label{prop:rigidity}
Any solution $\Phi$ of Problem~\ref{prob:char} is determined by its
values on operational graphs, is invariant under all cost-preserving
system isomorphisms, and satisfies $\Phi \le f\circ\E$ for
$f(t)=\sup\{\Phi_{I_t}(x,y)\}$, where $I_t$ is the two-state system
with a single operation of cost $t$. In particular $\Phi$ vanishes on
the diagonal and is infinite only where $\E$ is infinite, provided
$f$ is finite on finite arguments.
\end{proposition}

\begin{proof}[Proof sketch]
Isomorphisms are invertible admissibility morphisms whose inverses
are also admissibility morphisms, so (M) applied both ways gives
invariance. For the bound: given states $C,C'$ with
$\E(C,C')=t<\infty$, any operational geodesic word induces an
admissibility morphism from a chain system into $(\C,\Ops,c)$ hitting
$C,C'$, and a morphism from the chain onto $I_t$; monotonicity along
both pins $\Phi_{\C}(C,C')$ against $\Phi_{I_t}$. Where
$\E(C,C')=\infty$ no such comparison exists and $\Phi$ is
unconstrained from above.
\end{proof}

The gap between Propositions~\ref{prop:cone} and \ref{prop:rigidity}
is the substance of the problem: whether monotonicity plus invariance
force $\Phi$ to factor through $\E$ \emph{exactly}. On the
twelve-state system this can be settled by exhaustive computation, and
the answer is negative.

\begin{theorem}[Endomorphism constraints do not force factorization]
\label{thm:nofactor}
The system of Section~1 admits exactly $22$ admissibility
endomorphisms, of which the identity is the only $\E$-sufficient one.
The constraint digraph they generate on the $132$ ordered state pairs
has $132$ singleton equality classes, and there exists an explicit
functional --- the chain depth of a pair in the condensation order of
the constraints --- that satisfies (M) and (I) with respect to every
endomorphism yet takes five distinct values on the level set
$\{\E=2\}$ alone.
\end{theorem}

\begin{proof}
Computational; script \texttt{characterize.py} in the appendix.
Cost-nonincrease bounds each generator image to a word of cost at
most the generator's cost, and reachability of every state from
$(0,0)$ makes an endomorphism determined by the triple
$\bigl(F(0,0),F_{*}(a),F_{*}(b)\bigr)$; the $96$ candidates are
checked for loop closure, yielding $22$. Sufficiency, the constraint
digraph, its strongly connected components, and the verification of
(M) and (I) for the chain-depth functional are then finite checks.
\end{proof}

The separating witnesses are structurally meaningful. Within
$\{\E=2\}$ the functional distinguishes a two-step phase adjustment
$(0,0)\to(0,2)$ from a single promotion $(0,0)\to(1,0)$ and from a
promotion adjacent to the singular stratum $(1,0)\to(2,0)$: it is
sensitive to the \emph{operational type} of a geodesic, precisely the
data that $\E$ forgets and the commutator defect sees.

Theorem~\ref{thm:nofactor} uses only endomorphisms; Problem
\ref{prob:char} quantifies over all morphisms between all systems. Two
further results show that enlarging the constraint class does not
rescue factorization.

\begin{proposition}[Factorization of forced constraints]
\label{prop:factor}
Let $X$ be a finite continuation system. Any chain of monotonicity
constraints leading from an ordered pair of $X$ through pairs of other
systems and back to an ordered pair of $X$ is implied by the
constraint of a single endomorphism of $X$. Consequently the
equalities forced on pairs of $X$ by naturality over the entire
category coincide with those forced by $\operatorname{End}(X)$; for
the twelve-state system, none.
\end{proposition}

\begin{proof}
Each constraint edge is induced by a morphism applied to a pair, and
admissibility morphisms compose (intertwining assignments compose
word-wise; cost-nonincrease is transitive). A constraint chain
$X \to Y_1 \to \cdots \to X$ is therefore induced by the composite
morphism, an endomorphism of $X$. The claim about the model is
Theorem~\ref{thm:nofactor}.
\end{proof}

\begin{theorem}[A natural separating family: co-tests]
\label{thm:cotest}
For a finite continuation system $T$ with a distinguished ordered pair
$(t,t')$, define on every system $X$
\[
\Phi^{T,t,t'}_{X}(x,y) =
\begin{cases}
0 & \text{if some admissibility morphism } F\colon T\to X
\text{ has } F(t)=x,\ F(t')=y,\\
1 & \text{otherwise.}
\end{cases}
\]
Then $\Phi^{T,t,t'}$ satisfies (M) over the entire category. On the
twelve-state system, taking $T$ to be the system itself: the co-test
based at $\bigl((0,0),(1,0)\bigr)$ takes the value $0$ exactly on the
promotion-type and saturation-collapse pairs and $1$ on all
phase-type pairs, separating the level sets $\{\E=1\}$ and
$\{\E=2\}$; it satisfies (I) with respect to every $\E$-sufficient
endomorphism; and the $132$ co-tests based at the model's own pairs
are pairwise distinct as functionals on it, so the natural monotone
invariants of the category separate \emph{every} ordered pair of the
model.
\end{theorem}

\begin{proof}
Monotonicity: if $F\colon X\to Y$ and $\Phi_{X}(x,y)=0$ via
$H\colon T\to X$, then $F\circ H$ witnesses
$\Phi_{Y}\bigl(F(x),F(y)\bigr)=0$; the case $\Phi_{X}=1$ is vacuous.
The statements about the model are finite checks
(\texttt{crosstest.py}): morphisms $T=\C\to\C$ are the $22$
endomorphisms, so all values are computable.
\end{proof}

The co-tests show that monotonicity alone separates every ordered
pair: the natural (M)-invariants of the category are strictly finer
than the enrichment and detect operational data --- promotion versus
phase, proximity to the singular stratum --- that $\E$ forgets. The
question left open by Theorems~\ref{thm:nofactor} and
\ref{thm:cotest} was whether condition (I), over $\E$-sufficient
morphisms \emph{between} systems, cuts this profusion back. It does
--- completely --- and the mechanism is that sufficiency, scarce
inside any one irreversible system, can be \emph{manufactured}
between systems by gluing.

\begin{lemma}[Sufficiency is isometric embedding]\label{lem:embed}
Every $\E$-sufficient admissibility morphism is injective and
isometric for $\E$. (Costs are strictly positive, so
$\E(C,C')=0$ iff $C=C'$.)
\end{lemma}

\begin{lemma}[Amalgamation]\label{lem:amalgam}
Let $X,X'$ be finite continuation systems with ordered pairs $(x,y)$,
$(x',y')$ such that $\E_{X}(x,y)=\E_{X'}(x',y')$ and
$\E_{X}(y,x)=\E_{X'}(y',x')$. Let $Y = X \sqcup X' / (x\!\sim\!x',\,
y\!\sim\!y')$, where the generator sets $\Ops_X,\Ops_{X'}$ are taken
as disjoint copies before forming the quotient — so that two
generators of the same name or the same action on states are still
treated as distinct generators of $Y$ unless explicitly identified —
with the union of these disjoint copies as the generator set of $Y$. Then both
inclusions $X\hookrightarrow Y$ and $X'\hookrightarrow Y$ are
$\E$-sufficient admissibility morphisms. If instead only the
inequalities $\E_{X'}(x',y') \ge \E_{X}(x,y)$ and
$\E_{X'}(y',x') \ge \E_{X}(y,x)$ hold, the inclusion of $X$ alone is
$\E$-sufficient (the one-sided form).
\end{lemma}

The proof requires knowing exactly how a $Y$-path can interleave the
two components. The next lemma settles this before any cost estimate
is made, so that the induction below has no hidden assumptions about
how many times a path may cross between $X$ and $X'$.

\begin{lemma}[Path decomposition across a two-point gluing]\label{lem:pathdecomp}
Write $g\in\{x,y\}$ for the two gate points, identified with
$\{x',y'\}$ in $Y$, and let $w=o_n\cdots o_1$ be any operational word
in $Y$ from $a$ to $b$ with $a,b\in\C_X$, acting on states
$C_0=a,C_1,\dots,C_n=b$. Since $\Ops_X$ and $\Ops_{X'}$ are disjoint
(the generator sets are defined on the disjoint union before
gluing), each step $o_i$ lies in exactly one of the two sets; grouping
consecutive steps of the same type partitions $\{1,\dots,n\}$ into an
alternating sequence of maximal \emph{$X$-segments} and
\emph{excursions}. Every excursion begins and ends at a gate point.
\end{lemma}

\begin{proof}
Consider a maximal run of $\Ops_{X'}$-steps at positions
$i,\dots,j$. If $i>1$, the preceding step $o_{i-1}\in\Ops_X$ places
$C_{i-1}\in\C_X$; if $i=1$, $C_{i-1}=a\in\C_X$ by hypothesis. Either
way $C_{i-1}\in\C_X$. But $o_i\in\Ops_{X'}$ requires
$C_{i-1}\in\C_{X'}$, its domain. Since $\C_X\cap\C_{X'}=\{x,y\}$ after
gluing, $C_{i-1}\in\{x,y\}$. Symmetrically, $C_j$ is the codomain of
$o_j\in\Ops_{X'}$, so $C_j\in\C_{X'}$; if $j<n$ the following step
$o_{j+1}\in\Ops_X$ forces $C_j\in\C_X$, and if $j=n$ then $C_j=b\in\C_X$
by hypothesis. Either way $C_j\in\{x,y\}$. Hence every excursion runs
from a gate point to a gate point, regardless of how many excursions
$w$ contains or in what order they occur.
\end{proof}

\begin{proof}[Proof of Lemma~\ref{lem:amalgam}, two-sided case]
Fix $a,b\in\C_X$ and let $w$ be any operational word in $Y$ from $a$
to $b$. By Lemma~\ref{lem:pathdecomp}, $w$ decomposes into $X$-segments
separated by $k\ge0$ excursions $e_1,\dots,e_k$, each running between
gate points. We argue by induction on $k$ that
$\operatorname{cost}(w)\ge\E_X(a,b)$.

\emph{Base case $k=0$.} Then $w$ is entirely an $X$-word from $a$ to
$b$, so $\operatorname{cost}(w)\ge\E_X(a,b)$ directly from
Definition~\ref{def:free}.

\emph{Inductive step.} Let $e_1$ run from gate $g$ to gate $g'$
(each in $\{x,y\}$), with cost $c(e_1)\ge \E_{X'}(g,g')$, viewing
$g,g'$ as the corresponding points of $X'$. By hypothesis
$\E_{X'}(x',y')=\E_X(x,y)$, $\E_{X'}(y',x')=\E_X(y,x)$, and trivially
$\E_{X'}(x',x')=\E_X(x,x)=0=\E_{X'}(y',y')=\E_X(y,y)$; in every case
$\E_{X'}(g,g')=\E_X(g,g')$. Since $\C$ is finite and every cost is
strictly positive, an infimizing $X$-word from $g$ to $g'$ may be
taken simple (non-repeating), hence exists with cost exactly
$\E_X(g,g')\le c(e_1)$; call it $u_1$. Replace $e_1$ by $u_1$ in $w$
to obtain $w'$: the states immediately before and after $e_1$ are $g$
and $g'$, unchanged by the substitution, so $w'$ is a well-formed
word from $a$ to $b$ with $k-1$ excursions and
$\operatorname{cost}(w')\le\operatorname{cost}(w)$. By the inductive
hypothesis, $\operatorname{cost}(w')\ge\E_X(a,b)$, so
$\operatorname{cost}(w)\ge\operatorname{cost}(w')\ge\E_X(a,b)$.

Since $w$ was arbitrary, $\E_Y(a,b)=\inf_w\operatorname{cost}(w)\ge\E_X(a,b)$.
The reverse inequality $\E_Y(a,b)\le\E_X(a,b)$ holds because every
$X$-word is a $Y$-word. Hence $\E_Y(a,b)=\E_X(a,b)$ for \emph{every}
pair $a,b\in\C_X$, not merely $(x,y)$: the inclusion $X\hookrightarrow
Y$ is $\E$-sufficient. Exchanging the roles of $X$ and $X'$ gives the
same conclusion for $X'\hookrightarrow Y$.
\end{proof}

\begin{proof}[Proof of Lemma~\ref{lem:amalgam}, one-sided case]
The induction is unchanged for words with endpoints in $X$: the
excursion bound $c(e_\ell)\ge\E_{X'}(g_\ell,g_\ell')\ge\E_X(g_\ell,g_\ell')$
still licenses replacing $e_\ell$ by an $X$-word of no greater cost,
so $X\hookrightarrow Y$ remains $\E$-sufficient. The symmetric
statement for $X'$ need not hold: if $X$ offers a strictly cheaper
route between the gate points than $X'$ does internally, an
$X'$-pair may acquire a shortcut through $X$, and $X'\hookrightarrow
Y$ can fail to be sufficient.
\end{proof}

\begin{theorem}[Characterization by the bidirected cost]
\label{thm:main}
An assignment $\Phi$, giving each finite continuation system a
functional on ordered pairs with $\Phi(x,x)=0$, satisfies (M) and (I)
over the category of finite continuation systems if and only if there
is a function $f$ on cost pairs, nondecreasing in each argument with
$f(0,0)=0$, such that
\[
\Phi_{X}(x,y) \;=\; f\bigl(\E_{X}(x,y),\,\E_{X}(y,x)\bigr)
\]
for every system $X$ and every pair.
\end{theorem}

\begin{proof}
\emph{Sufficiency of the form.} Lemma~\ref{lem:nonexp} applied to
both orderings gives (M) for nondecreasing $f$; an $\E$-sufficient
morphism preserves both coordinates, giving (I); pairs collapsing to
the diagonal give $\Phi=0\le$ anything.

\emph{Necessity.} Let $V_{s,t}$ denote the two-state system with a
forward generator of cost $s$ and a reverse generator of cost $t$
(each omitted when infinite), so $\E_{V}(v_0,v_1)=s$,
$\E_{V}(v_1,v_0)=t$. Given any pair $(x,y)$ in any $X$ with cost pair
$(s,t)$, amalgamate $X$ with $V_{s,t}$ along
$(x,y)\sim(v_0,v_1)$; by Lemma~\ref{lem:amalgam} both inclusions are
sufficient, so (I) twice gives
$\Phi_{X}(x,y) = \Phi_{Y}(\text{glued pair}) = \Phi_{V_{s,t}}(v_0,v_1)
=: f(s,t)$. Thus $\Phi$ factors through the cost pair. Monotonicity
of $f$: for $s'\le s$, $t'\le t$ the identity on states with
generator reassignment defines an admissibility morphism
$V_{s,t}\to V_{s',t'}$ (a generator maps to one of no greater cost,
or to a generator absent in the source there is nothing to assign),
and (M) gives $f(s',t')\le f(s,t)$.
\end{proof}

\begin{corollary}[Fate of the co-tests]\label{cor:cotest}
No co-test invariant that fails to factor through the bidirected cost
satisfies (I) over the category. Explicitly, for
$\Phi^{\C,(0,0),(1,0)}$: gluing a second copy of the model at
$\bigl((0,0),(1,0)\bigr)$ onto the pair $\bigl((0,0),(0,2)\bigr)$ ---
boundary costs $(2,\infty)$ against $(2,2)$, so the base inclusion is
sufficient by the one-sided form of Lemma~\ref{lem:amalgam} --- makes
the second copy's inclusion a cover of the image pair, forcing the
value $0$ where the base system has $1$ (verified in
\texttt{amalgam.py}).
\end{corollary}

Theorem~\ref{thm:main} is a finite, continuation-theoretic analogue
of the \emph{structure} of Chentsov's characterization argument, with
two departures that keep the analogy from running deeper than it
should. What is proved here is a characterization under a chosen
notion of admissibility morphism and a specifically constructed
notion of sufficiency, on a finite state space; it is Chentsov-like
in shape, not a finite instance of the same depth of result, since
the classical theorem characterizes the Fisher metric among all
smooth monotone metrics on an infinite-dimensional statistical
manifold, a generality this paper does not attempt. First, the role
that the \emph{abundance} of sufficient statistics plays in the
classical proof is played here by a \emph{construction}: sufficient
morphisms are rare inside irreversible systems
(Theorem~\ref{thm:nofactor}) but can always be manufactured between
systems by amalgamation, and the characterization rests entirely on
the manufactured ones. Second, the invariant characterized is not a
symmetric metric but the \emph{bidirected} executable cost --- the
ordered pair of forward and reverse repair cost. Directedness
survives characterization;
operational-bracket data does not. The torsion-sensitive invariants of
Theorem~\ref{thm:cotest} are thereby given an exact status: they are
natural under monotone comparison but not under lossless recoding.
What sufficiency-invariance preserves is what survives every
distinction-preserving change of description, and the theorem says
this is precisely two numbers per ordered pair: how much it costs to
get there, and how much it costs to get back. The two-primitive
architecture accordingly resolves as a stratification rather than a
theorem or a refutation: the (M)-natural theory sees the pair
(enrichment, operational structure); the (M)+(I)-natural theory sees
the bidirected enrichment alone. The precedent of Markov categories
\cite{fritz} suggests how the argument should be organized
synthetically; the infinite-dimensional Chentsov theorem \cite{ajls}
remains the caution for the continuous case.

\section{Smooth Limits and Dual Geometry}\label{sec:smooth}

The enriched theory is native to the finite, procedural setting. The
differential-geometric apparatus enters only through a relaxation, and
we present it as such: a tangent theory of the regular stratum, not a
foundation.

\subsection{The simplex relaxation}

Embed $\C \hookrightarrow \Delta(\C)$, the simplex of probability
weights on states, sending $C$ to the vertex $\delta_C$. Relaxed
states are mixtures $\pi\in\Delta(\C)$; generators extend to affine
maps $o_{*}\pi = \pi\circ o^{-1}$ (pushforward), and a directed
contrast extends to a smooth two-point function
$\mathcal{R}\colon \Delta(\C)\times\Delta(\C)\to[0,\infty]$ agreeing
with $\R$ on vertices. This is the same move Chentsov makes: the
finite object is primary, the geometry lives on distributions over it.

\subsection{Induced dual geometry}

Where $\mathcal{R}$ is smooth with nondegenerate diagonal Hessian, the
Eguchi construction \cite{amari} yields, in local coordinates on
$\Delta(\C)$,
\[
g_{ab} = -\left.\frac{\partial^{2}\mathcal{R}(\pi,\pi')}
{\partial \pi^{a}\,\partial \pi'^{b}}\right|_{\pi'=\pi},
\qquad
\Gamma_{abc} = -\left.\frac{\partial^{3}\mathcal{R}(\pi,\pi')}
{\partial \pi^{a}\partial \pi^{b}\,\partial \pi'^{c}}\right|_{\pi'=\pi},
\]
with the dual $\Gamma^{*}$ obtained by exchanging primed and unprimed
arguments. The metric is symmetric regardless of the asymmetry of
$\mathcal{R}$; the directedness of repair survives in the cubic
tensor $T = \Gamma^{*}-\Gamma$ and the dual pair
$(\nabla,\nabla^{*})$, jointly metric-compatible:
$X\,g(Y,Z) = g(\nabla_X Y, Z) + g(Y,\nabla^{*}_X Z)$. Both induced
connections are torsion-free by construction. Directed repair
asymmetry and procedural order-dependence are therefore distinct
departures from Riemannian geometry, and only the first is visible to
the contrast function.

\subsection{Torsion as representation defect}

Order-dependence enters through the operational algebra. Suppose the
relaxed generators near the identity are generated by vector fields
$X_{a}$ on $\Delta(\C)$, with operational bracket
$[X_{a},X_{b}]_{\Ops}$ determined by the composition law. For any
candidate connection $\nabla$, define
\[
\Tor(X_{a},X_{b})
= \nabla_{X_{a}}X_{b} - \nabla_{X_{b}}X_{a} - [X_{a},X_{b}]_{\Ops}.
\]
Torsion so defined does not record noncommutativity as such; it
records the failure of geometric transport to represent the
operational bracket. Three cases separate: commuting operations with
torsion-free geometry; noncommuting operations exactly represented,
again torsion-free; and genuine mismatch. The twelve-state system
suggests where mismatch lives:

\begin{conjecture}[Singular support]\label{conj:support}
$\operatorname{supp}(\Tor)$ is contained in the operational singular
set --- the closure of the locus where some generator fails to be
injective or applicable.
\end{conjecture}

In the model this is the saturation stratum $\{\ell=2\}$, on which the
commutator defect of $(a,b)$ is entirely concentrated.

\subsection{Decomposing the projection defect}

Section~1 exhibited both behaviours of the defect $\Delta_{\A}$:
identically zero for a level-aligned admissible set, strictly negative
on a diagonal one. In the smooth regime the generalized Pythagorean
theorem holds when the ambient geometry is dually flat \emph{and} the
admissible set is autoparallel in the relevant dual connection; a
nonzero defect therefore has two independent sources, ambient
non-flatness and submanifold curvature, plus their interaction. Any
diagnostic use of $\Delta_{\A}$ must decompose it accordingly. The
enriched regime adds the sign restriction of
Proposition~\ref{prop:sign} on top: for triangular $\mathcal{R}$ the
entire positive half of the phenomenon is excluded, so dually flat
positive decompositions certify that the relaxed contrast has left the
Lawvere regime. The boundary between the regimes is crossed exactly
when the relaxation destroys the triangle inequality --- as it does,
for instance, when $\mathcal{R}$ is taken to be a Bregman or
Kullback--Leibler-type divergence extending a vertex-triangular $\R$.
Which extensions preserve triangularity, and what geometry is possible
under that constraint, appears to be open.

\subsection{The continuous problem}

\begin{conjecture}[Operational characterization, smooth case]
For a suitable class of smooth continuation systems, the tuple
$(g,\nabla,\nabla^{*},\Tor)$ --- with $(g,\nabla,\nabla^{*})$ induced
by the regular part of the contrast and $\Tor$ the representation
defect of the operational bracket --- is, up to normalization, the
unique local geometry monotone under all admissibility morphisms and
invariant under the repair-sufficient subclass.
\end{conjecture}

We state this as a target, not a claim, but
Theorem~\ref{thm:main} now constrains its form. In the finite case,
sufficiency-invariance characterized the bidirected cost, not the
operational bracket: the smooth analogue should therefore expect the
germ of the bidirected cost along the diagonal to be the
characterized object, with the metric $g$ as its symmetric
second-order jet and the cubic tensor $T$ carried by the third-order
asymmetry --- the dual geometry reappearing as the jet expansion of
the characterized invariant rather than as an independently selected
structure --- while $\Tor$, like the co-tests, is expected to be
natural only under monotone comparison. The amalgamation technique
should transfer: gluing smooth systems along cost-matched pairs is a
boundary construction, and its regularity is where the difficulty
will concentrate. The infinite-dimensional Chentsov theorem required
four decades and carefully engineered regularity \cite{ajls}; there
is no reason to expect the continuation version to be easier. What
this paper establishes is finite and secure: that the discrete
theory is enriched-categorical rather than differential-geometric in
kind; that the enriched and divergence regimes are separated by a
computable sign; that coherence between comparison and execution is a
local certificate in the first regime and a global axiom in the
second; and that the sufficiency-invariant content of a finite
continuation system is exactly its bidirected executable cost, no
more and no less.

\section{Interpretations within the Admissibility Program}\label{sec:interp}

The finite theory just established has consequences beyond the
category of continuation systems in which it was proved. This section
draws out five of them, in each case deriving a statement from
Theorem~\ref{thm:main} or Proposition~\ref{prop:sign} rather than
merely asserting a thematic resemblance.

\subsection{Admissibility and reachability}\label{subsec:admissibility}

For a state $x\in\C$, define the \emph{forward continuation
neighborhood} $V(x)=\{y\in\C:\E(x,y)<\infty\}$, the set of states
reachable from $x$ at finite executable cost, and dually the
\emph{backward} neighborhood $V^{-1}(x)=\{y\in\C:\E(y,x)<\infty\}$.
Admissibility of a continuation is then not a property of a
destination state in isolation but of the neighborhood it occupies:
two states $x,x'$ with $V(x)=V(x')$ and $V^{-1}(x)=V^{-1}(x')$ are
interchangeable with respect to what can follow and what could have
preceded them, even if $x\ne x'$ and $\R(x,x')>0$.

Given a comparison $d$ on subsets of $\C$, define the
\emph{admissibility distortion}
$D(x,x')=d(V(x),V(x'))+d(V^{-1}(x),V^{-1}(x'))$. This quantity is a
functional of $\{\E(x,\cdot),\E(\cdot,x)\}$ alone, since $V(x)$ and
$V^{-1}(x)$ are level sets of $\E$. Theorem~\ref{thm:main} therefore
applies to it directly: an admissibility distortion built from
reachability sets in this way is automatically expressible as a
function of sufficiency-invariant content, because it is a function
of the bidirected cost. An admissibility distortion built instead from
$\R$-neighborhoods $\{y:\R(x,y)<\theta\}$ carries no such guarantee,
by Remark~\ref{rem:specialize}: it may distinguish states that every
$\E$-sufficient redescription identifies. The bidirected cost is, in
this sense, exactly the resolution at which reachability-based
admissibility judgments remain invariant, and no finer.

This bears on the reconstruction literature in one further respect.
A state $x$ is reconstructible from a later state $y$ exactly when
$x\in V^{-1}(y)$; the theorem's content, applied here, is that what
survives redescription is not the identity of $x$ but the pattern of
which states lie in $V^{-1}(y)$ and at what bidirected cost.
Reconstruction, on this reading, is a question about continuation
structure rather than about recovering a state's identity as such.

The witness adequacy functional of the smooth, continuous-trajectory
setting --- $\Lambda(w)=\sup_{H\in\A(w)}d(H,\hat H(w))$, the
worst-case distance from a witness's admissible continuation set to a
model's output --- is the natural smooth-case cousin of the
reachability neighborhoods $V(x)$ and $V^{-1}(x)$ defined here: both
quantify how tightly partial information constrains a continuation
set, one by a sup over a metric and the other by a level set of
$\E$. Making that correspondence precise is part of the smooth
characterization problem of Section~\ref{sec:smooth} and is not
attempted here.

\subsection{Projection defect and CLIO}\label{subsec:clio}

The projection-defect construction of Section~1 admits a reading in
the vocabulary of CLIO, where a system's \emph{substrate} is the
space in which change occurs and an \emph{interface} is a subspace
onto which the substrate is represented. Under that correspondence,
$\C$ with $\E$ is the substrate, an admissible subset $\A$ is an
interface, and the projection $C^{*}=\arg\min_{t\in\A}\R(C,t)$ is the
representation of $C$ at that interface. The defect
$\Delta_{\A}(C,C')=\R(C,C')-\R(C,C^{*})-\R(C^{*},C')$ measures
representational distortion: the amount by which comparing two
substrate points directly disagrees with comparing them through the
interface.

Proposition~\ref{prop:sign} restricts which direction that distortion
can run. If the substrate's comparison structure is enriched,
$\Delta_{\A}\le0$ for every interface: representation through an
interface can only make two points look \emph{more} alike than direct
comparison does, never less. Positive projection interference — an
interface making two points look less alike than they are, the
signature phenomenon of dually flat information geometry — requires
$\R$ to fail the triangle inequality, i.e.\ requires leaving the
executable substrate for a divergence-type contrast that is not
itself a hom-structure of any enrichment. The enriched/divergence
fork of Section~1 is, under this correspondence, a fork in what kinds
of representational distortion an interface can produce: monotone-only
distortion in the enriched regime, and reversal of apparent similarity
only in the divergence regime, where coherence between substrate and
interface becomes a global postulate rather than a locally checkable
one (Lemma~\ref{lem:coherence} has no divergence-regime analogue).

\subsection{Distinction geometry and repair}\label{subsec:distinction}

\begin{definition}[Operational faithfulness]\label{def:faithful}
Say $\R$ is \emph{operationally faithful} on the ordered pair $(x,y)$
if some operational word $w$ from $x$ to $y$ satisfies
$\operatorname{cost}(w)=\R(x,y)$ --- that is, if the recognition cost
$\R(x,y)$ is actually achieved by an executable path, not merely
bounded by one.
\end{definition}

\begin{proposition}\label{prop:faithful}
If $\R$ is a generator-dominated enrichment, then $\G(x,y)=0$ if and
only if $\R$ is operationally faithful on $(x,y)$.
\end{proposition}

\begin{proof}
If some word $w$ realizes $\R(x,y)$, then
$\E(x,y)\le\operatorname{cost}(w)=\R(x,y)\le\E(x,y)$, the second
inequality by Lemma~\ref{lem:coherence}, so
$\G(x,y)=\E(x,y)-\R(x,y)=0$. Conversely, if $\G(x,y)=0$ then
$\E(x,y)=\R(x,y)$; since $\C$ is finite and every cost is strictly
positive, the infimum defining $\E(x,y)$ is attained by some word
$w$ (Lemma~\ref{lem:amalgam}'s proof makes the same finiteness
observation), and $\operatorname{cost}(w)=\E(x,y)=\R(x,y)$, so $w$
realizes $\R(x,y)$ and $\R$ is operationally faithful on $(x,y)$.
\end{proof}

The gap $\G=\E-\R$ has, to this point, been treated as a computed
diagnostic; Proposition~\ref{prop:faithful} gives its vanishing a
name --- recognition and execution agree on a pair exactly when $\R$
is operationally faithful there --- and turns "recognition agrees
with execution," used informally throughout the paper, into a
condition with a proof rather than a figure of speech. It bears directly on a question
Distinction Geometry and Repair Theory pose separately: the first asks
how distinguishable two states are, the second how recoverable one is
from the other. The twelve-state system shows these are not one
structure under two names. Forty-eight of the $132$ ordered pairs have
$\R<\infty$ but $\E=\infty$: every downward repair is recognizable,
often cheaply, yet inexecutable, since no operation lowers the level
$\ell$. States can therefore be maximally distinguishable in the
operational sense while being cheaply comparable under $\R$.
Conversely, among finite pairs, $\R((0,0),(2,3))=2.5$ against
$\E((0,0),(2,3))=7$: states a recognition process treats as nearly
identical are separated by substantial executable cost, because $\R$
exploits a cyclic symmetry the operational algebra does not respect.

$\G$ is best read, on this evidence, as a \emph{distinction--repair
mismatch}: the extent to which recognizability and recoverability
diverge, rather than a bookkeeping difference. Theorem~\ref{thm:main}
sharpens what the mismatch means for invariant theory. Recognizability,
as encoded in $\R$, need not be sufficiency-invariant at all; it sits
outside the characterization of Section~\ref{sec:char} except where it
coincides with $\E$ (Remark~\ref{rem:specialize}). What is
characterized is recoverability alone, in both directions. A theory
that identifies distinguishability with recoverability is, on this
evidence, conflating two independent structures; the present results
quantify the conflation without resolving it, and the joint invariant
theory of $(\R,\E)$ remains the natural next problem.

\subsection{Continuation geometry}\label{subsec:continuation}

The neighborhood $V(x)$ of Section~\ref{subsec:admissibility}
organizes into a continuation field $\Gamma$ assigning to each state
its reachable future, with the bidirected cost
$\bigl(\E(x,y),\E(y,x)\bigr)$ as local coordinates on pairs within it.
In this language, Theorem~\ref{thm:main} restates as: what survives
every admissibility-preserving redescription of a continuation system
is not the identity of its states but the geometry of $\Gamma$ ---
the reachable-future structure and its bidirected costs. This is a
restatement, not an extension, of the theorem, offered to align
vocabulary with related work where the same distinction is drawn
without the present finite-case proof.

\subsection{Further connections}\label{subsec:further}

Two remarks, offered as interpretation rather than as claims proved
here. Executable cost has, in RSVP, the shape of a path-integral
functional through capacity, transport, and entropy fields,
$\E(x,y)=\inf_\gamma\int_\gamma L(\Phi,v,S)\,dt$; the finite system of
this paper is then a discrete prototype of that reachability
structure, small enough to compute exhaustively, but nothing proved
here shows that RSVP's continuous reachability functionals satisfy an
analogue of Theorem~\ref{thm:main} --- that remains conjectural,
alongside the smooth case of Section~\ref{sec:smooth}. In systems with
multiple interacting repair modules, distinct modules may agree on
recognition ($\R$) while disagreeing on executable repair ($\E$); the
present results say that only the latter agreement is what a shared,
sufficiency-invariant repair policy can rely on, with $\G$ measuring
the gap between what a module recognizes as reparable and what it can
execute.

A caution belongs here. Where repair operators have been defined
elsewhere as contractions toward an admissible manifold, that
contraction property is an additional hypothesis, not a consequence
of anything proved in this paper: neither $\R$ nor $\E$ is assumed
here to shrink under iteration, and Theorem~\ref{thm:main} places no
requirement of that kind on either primitive. Any specialization of
the present framework to contraction-type repair operators should
state the contraction hypothesis explicitly rather than inherit it by
association.

\section{What Survives}

Theorem~\ref{thm:main} admits two readings, and they point the
program in different directions.

The first is deflationary and should be stated plainly. If one asks
what a continuation system \emph{is}, invariantly --- what remains
when every distinction-preserving redescription has been quotiented
away --- the answer is austere: a bidirected cost structure, the
forward and reverse price of every transition. The rich operational
texture that the monotone invariants detect --- the difference
between a promotion and a phase adjustment of equal cost, the
concentration of noncommutativity on a singular stratum, the entire
torsion-like layer of structure --- none of it is
sufficiency-invariant. It is real, and it is natural in the weaker
sense, but it is perspective-bound: visible to comparison, erased by
lossless recoding. A theory that wants to speak invariantly about
continuation systems has exactly two numbers per ordered pair to
speak with.

The second reading is the one the Admissibility Program should take.
The theorem does not say the operational stratum is illusory; it says
the operational stratum is precisely the content of a
\emph{description}, over and above the content of the system
described. The gap between the (M)-natural theory and the
(M)+(I)-natural theory measures how much structure lives in the
choice of presentation --- in which operations are taken as
generators, where their singularities fall, how their order matters
--- and the theorem locates that structure exactly rather than
gesturing at it. In the vocabulary of the surrounding work: the
bidirected cost is what is \emph{restorable} from any faithful
witness of the system; the bracket data is what a witness
additionally knows in virtue of how it witnessed. That the invariant
content of a continuation system reduces to reachability cost in
both directions --- not to a truth-structure, not to a
state-structure, but to what can be reached and what can be
recovered --- is, in the finite case exactly, the admissibility
thesis in the form of a theorem.

The question that opened this inquiry --- whether an
admissibility-theoretic framework could derive the covariant
derivative --- can now be answered with the precision it originally
lacked. In the weak sense, yes, and trivially: any transport of
continuations across varying contexts has the shape
$\partial + \Omega$. In the strong sense the question was wrongly
posed. The covariant derivative of Riemannian geometry is the unique
transport compatible with a symmetric quadratic comparison
structure; the present setting has a directed, frequently
non-quadratic, frequently non-smooth comparison structure, and what
is unique given \emph{its} invariances is not a connection but the
bidirected cost. Connection-like objects reappear only in the
tangent approximation, as jets of the characterized invariant, which
is where they belong. What the frameworks of this program derive
from first principles is not the $\Gamma$ of general relativity; it
is the thing of which $\Gamma$ is a smooth, symmetric, reversible
shadow.

Beyond the continuous conjecture of Section~4, three problems remain
open in the finite theory itself. The characterization concerns
functionals of a single ordered pair; the invariant theory of
$n$-point configurations, where amalgamation has less room to move,
is untouched, and it is there that bracket data may re-enter the
sufficiency-invariant world. The morphism class was fixed by one
notion of cost-nonincrease; coarser and finer classes ---
cost-scaling morphisms, morphisms preserving an admissible cone,
morphisms preserving a repair pairing --- will characterize
different invariants, and the map from morphism class to
characterized structure is the proper generalization of the selector
problem with which this inquiry began. And the repair contrast $\R$,
which the characterization deliberately quotiented out, returns as
data the moment one fixes a presentation: the joint invariant theory
of the pair $(\R,\E)$ --- of recognition against execution --- is
the natural next paper.

One sentence of the program survives this inquiry unchanged and
acquires a proof schema. Repair is not something performed inside a
geometry; repair is the operation from which the geometry is
inferred. What the finite theory adds is the terminus of the
inference. Pose it invariantly, and the geometry it infers is the
bidirected cost: what it costs to arrive, and what it costs to
return.

\section*{Appendix: Computation}

The complete computation is reproduced by the accompanying scripts: \texttt{toymodel.py} (Section~1: Dijkstra over the operational graph for $\E$; exhaustive evaluation of $\R$, $\G$, the commutator defect, and $\Delta_{\A}$ for both admissible sets), \texttt{characterize.py} (Theorem~\ref{thm:nofactor}), \texttt{crosstest.py} (Theorem~\ref{thm:cotest}), and \texttt{amalgam.py} (Lemma~\ref{lem:amalgam} instances and Corollary~\ref{cor:cotest}). The full matrices follow; the diagonal is zero.

\begin{table}[h]\centering\scriptsize
\setlength{\tabcolsep}{3pt}
\begin{tabular}{l|cccccccccccc}
 & $00$ & $01$ & $02$ & $03$ & $10$ & $11$ & $12$ & $13$ & $20$ & $21$ & $22$ & $23$ \\ \hline
$00$ & 0 & 1 & 2 & 3 & 2 & 3 & 4 & 5 & 4 & 5 & 6 & 7 \\
$01$ & 3 & 0 & 1 & 2 & 5 & 2 & 3 & 4 & 6 & 4 & 5 & 6 \\
$02$ & 2 & 3 & 0 & 1 & 4 & 5 & 2 & 3 & 6 & 7 & 4 & 5 \\
$03$ & 1 & 2 & 3 & 0 & 3 & 4 & 5 & 2 & 5 & 6 & 7 & 4 \\
$10$ & $\infty$ & $\infty$ & $\infty$ & $\infty$ & 0 & 1 & 2 & 3 & 2 & 3 & 4 & 5 \\
$11$ & $\infty$ & $\infty$ & $\infty$ & $\infty$ & 3 & 0 & 1 & 2 & 4 & 2 & 3 & 4 \\
$12$ & $\infty$ & $\infty$ & $\infty$ & $\infty$ & 2 & 3 & 0 & 1 & 4 & 5 & 2 & 3 \\
$13$ & $\infty$ & $\infty$ & $\infty$ & $\infty$ & 1 & 2 & 3 & 0 & 3 & 4 & 5 & 2 \\
$20$ & $\infty$ & $\infty$ & $\infty$ & $\infty$ & $\infty$ & $\infty$ & $\infty$ & $\infty$ & 0 & 1 & 2 & 3 \\
$21$ & $\infty$ & $\infty$ & $\infty$ & $\infty$ & $\infty$ & $\infty$ & $\infty$ & $\infty$ & 2 & 0 & 1 & 2 \\
$22$ & $\infty$ & $\infty$ & $\infty$ & $\infty$ & $\infty$ & $\infty$ & $\infty$ & $\infty$ & 2 & 3 & 0 & 1 \\
$23$ & $\infty$ & $\infty$ & $\infty$ & $\infty$ & $\infty$ & $\infty$ & $\infty$ & $\infty$ & 1 & 2 & 3 & 0 \\
\end{tabular}
\caption{Executable cost $\E$. Rows: source state $\ell p$; columns: target.}
\end{table}

\begin{table}[h]\centering\scriptsize
\setlength{\tabcolsep}{3pt}
\begin{tabular}{l|cccccccccccc}
 & $00$ & $01$ & $02$ & $03$ & $10$ & $11$ & $12$ & $13$ & $20$ & $21$ & $22$ & $23$ \\ \hline
$00$ & 0 & 0.5 & 1 & 0.5 & 1 & 1.5 & 2 & 1.5 & 2 & 2.5 & 3 & 2.5 \\
$01$ & 0.5 & 0 & 0.5 & 1 & 1.5 & 1 & 1.5 & 2 & 2.5 & 2 & 2.5 & 3 \\
$02$ & 1 & 0.5 & 0 & 0.5 & 2 & 1.5 & 1 & 1.5 & 3 & 2.5 & 2 & 2.5 \\
$03$ & 0.5 & 1 & 0.5 & 0 & 1.5 & 2 & 1.5 & 1 & 2.5 & 3 & 2.5 & 2 \\
$10$ & 5 & 5.5 & 6 & 5.5 & 0 & 0.5 & 1 & 0.5 & 1 & 1.5 & 2 & 1.5 \\
$11$ & 5.5 & 5 & 5.5 & 6 & 0.5 & 0 & 0.5 & 1 & 1.5 & 1 & 1.5 & 2 \\
$12$ & 6 & 5.5 & 5 & 5.5 & 1 & 0.5 & 0 & 0.5 & 2 & 1.5 & 1 & 1.5 \\
$13$ & 5.5 & 6 & 5.5 & 5 & 0.5 & 1 & 0.5 & 0 & 1.5 & 2 & 1.5 & 1 \\
$20$ & 10 & 10.5 & 11 & 10.5 & 5 & 5.5 & 6 & 5.5 & 0 & 0.5 & 1 & 0.5 \\
$21$ & 10.5 & 10 & 10.5 & 11 & 5.5 & 5 & 5.5 & 6 & 0.5 & 0 & 0.5 & 1 \\
$22$ & 11 & 10.5 & 10 & 10.5 & 6 & 5.5 & 5 & 5.5 & 1 & 0.5 & 0 & 0.5 \\
$23$ & 10.5 & 11 & 10.5 & 10 & 5.5 & 6 & 5.5 & 5 & 0.5 & 1 & 0.5 & 0 \\
\end{tabular}
\caption{Repair contrast $\R$. Rows: source state $\ell p$; columns: target.}
\end{table}



\begin{thebibliography}{9}
\bibitem{lawvere} F.~W. Lawvere, \emph{Metric spaces, generalized
logic, and closed categories}, Rend.\ Sem.\ Mat.\ Fis.\ Milano 43
(1973).
\bibitem{chentsov} N.~N. Chentsov, \emph{Statistical Decision Rules
and Optimal Inference}, AMS Translations (1982).
\bibitem{amari} S.~Amari and H.~Nagaoka, \emph{Methods of Information
Geometry}, AMS/Oxford (2000).
\bibitem{fritz} T.~Fritz, \emph{A synthetic approach to Markov kernels,
conditional independence and theorems on sufficient statistics},
Adv.\ Math.\ 370 (2020).
\bibitem{ajls} N.~Ay, J.~Jost, H.~V. L\^e, L.~Schwachh\"ofer,
\emph{Information Geometry}, Springer (2017).
\end{thebibliography}

\end{document}
